%======================
% Chapter 1 : Introduction
%======================

\chapter{\scshape Introduction}

\section{Background}

Falls are a leading cause of injury-related mortality and morbidity among adults aged 65~and above worldwide. The WHO Global Report on Falls Prevention in Older Age (2007) identifies falls as the second leading cause of unintentional injury death globally, with approximately 684,000 fatal falls occurring each year\cite{who2007falls}. Non-fatal falls cause serious injuries including hip fractures, traumatic brain injuries, and long-term disability, placing a heavy burden on healthcare systems and caregivers.

Nepal is undergoing a rapid demographic shift. The Central Bureau of Statistics projects that Nepal's population aged 65 and above will double by 2035~\cite{cbs2021nepal}, driven by improving life expectancy and urbanization. Despite this growing demographic pressure, no domestic fall detection product exists in the Nepali market. Commercial solutions such as the Apple Watch and medical alert pendants are cost-prohibitive --- a single Apple Watch SE (2nd~generation) retails above NPR~30,000 --- and require proprietary ecosystems unavailable to most Nepali families. Furthermore, existing wearable datasets and commercial algorithms reflect younger Western cohorts and waist placements, failing to address the elderly kinematic gap (where frail elderly impact peaks range between $3.5\text{--}6.0$\,g compared to $8.0\text{--}12.0$\,g in young simulated falls). A local BLE link between the wearable and a household laptop gateway requires no internet connectivity at all, making SPARK deployable in care settings regardless of WiFi/ISP reliability.

Academic fall detection systems published in the literature typically rely on one of two paradigms: threshold-only detection on raw accelerometer data, which produces unacceptably high false-positive rates on activities of daily living such as sit-to-stand transitions and stair climbing; or cloud-dependent machine learning inference, which introduces unacceptable network dependencies for a safety-critical application. Even among edge AI deployments, existing works either execute un-gated, always-on neural networks that rapidly deplete wearable batteries, offload classification to a remote server, or rely on cost-prohibitive FPGA/SoC architectures. Moreover, no published open-source system provides per-event classifier explainability to clinicians and care staff.

SPARK addresses these gaps with a two-layer gated edge AI architecture running entirely on an ESP32-S3 microcontroller, a gradient-based SHAP explainability pipeline on a local companion gateway, and a standardized 34-activity Nepal-context IMU dataset collected at KEC.\@

\section{Problem Statement}

Falls are the leading cause of injury-related mortality among adults over 65~globally\cite{who2007falls}. Nepal's elderly population is projected to double by 2035~\cite{cbs2021nepal}, yet no affordable domestic fall detection solution exists. Commercial fall detection systems (Apple Watch, Samsung Galaxy Watch, Withings ScanWatch) are cost-prohibitive for Nepal's eldercare context, require proprietary ecosystems, and operate as closed black-box systems --- returning a binary fall verdict with no classifier explanation. Published academic systems use either threshold-only detection (unacceptably high false-positive rate) or cloud-dependent machine learning inference (network dependency unacceptable for safety-critical applications in Nepal's connectivity environment). 

Crucially, existing literature lacks an end-to-end framework combining:
\begin{enumerate}
\item \textbf{On-Device Two-Stage Gating on Commodity Hardware:} A continuously-sampled software magnitude-and-duration threshold gate combined with subsequent quantized INT8 1D CNN confirmation, executing entirely on a single sub-\$5 microcontroller (ESP32-S3) with zero classification offload.
\item \textbf{Local Companion-Device Explainability:} Per-event gradient feature attribution ($\text{Input} \times \nabla_{\text{Input}}$) computed locally on a companion gateway without cloud dependencies, rendering transparent biomechanical feature attributions.
\item \textbf{Kinematically-Calibrated Sensitivity:} Adaptive sensitivity profiles (High, Standard, Sport) and partial-freeze fine-tuning addressing the elderly impact attenuation gap.
\item \textbf{Automated Clinical Reporting \& Local Delivery:} Real-time generation of auditable one-page clinical PDF incident reports with zero-dependency REST/Web dashboard delivery validated on a South Asian cohort dataset.
\end{enumerate}

\section{Objectives}

This project focuses on building a wearable fall detection system that combines open hardware, embedded edge AI, and local companion explainability to automate real-time monitoring and caregiver alerting. The key objectives are:
\begin{enumerate}
\item To construct an ergonomic, wrist-worn wearable edge node based on the ESP32-S3 and MPU6050 with optimized DFM (TPU 95A enclosure, brass inserts).
\item To develop and deploy an autonomous two-layer gated AI architecture (ISR threshold pre-filter + 18.5\,KB INT8 CNN) for cloud-independent, real-time fall classification ($<200$\,ms latency).
\item To implement a local gateway explainability engine calculating gradient-based SHAP attributions and auto-generating clinical incident PDF reports across distinct fall archetypes.
\item To establish a standardized 34-activity motion dataset protocol (15 falls, 19 ADLs) with $N=12\text{--}20$ volunteers at Kathmandu Engineering College to validate model transferability.
\end{enumerate}

\section{Scope of the Project}\label{sec:scope}

The scope of our project includes:

\begin{itemize}
    \item \textbf{Continuous On-Device Data Acquisition:} Sampling 6-axis accelerometer and gyroscope data directly on the wearable node at a continuous 200\,Hz frequency over I$^2$C with a circular sliding-window ring buffer.
    \item \textbf{Dual-Layer Gated Edge Inference:} Layer~1 software threshold pre-filter ($|a| > 2.5$\,g for $>300$\,ms) running in the 200\,Hz ISR, gating Layer~2 quantized INT8 1D CNN inference ($18.5$\,KB footprint) only upon potential fall impact.
    \item \textbf{Zero-Cloud BLE Event Transport:} Asynchronous BLE GATT notification of confirmed fall JSON telemetry directly to a local companion laptop gateway.
    \item \textbf{Local Companion XAI Pipeline:} Rapid gradient-based feature attribution ($X \times \nabla_X Y$) executing on the gateway (benchmarked across Intel Core Ultra 7 CPU/iGPU/NPU engines via OpenVINO) to explain kinetic channel contributions.
    \item \textbf{Automated Clinical PDF Documentation:} Instant compilation of one-page clinical incident reports with biomechanical waveform attribution charts across 4 distinct fall archetypes (forward, lateral, syncope, rotational).
    \item \textbf{Local REST API \& Zero-Dependency Web Dashboard:} Standalone HTTP/CORS server on the companion gateway providing caregiver status views and report downloads.
    \item \textbf{Nepal-Context Experimental Validation:} Structured 34-activity data collection on gymnastic crash mats at KEC for empirical model evaluation and partial-freeze transfer learning.
\end{itemize}
\section{Application of the Project}
\begin{enumerate}

\item Assistive Technology: Provides real-time, hands-free fall monitoring to increase the independence and safety of vulnerable individuals living alone.

\item Medical \& Rehabilitation: Assists clinicians in tracking recovery progress and movement patterns during post-injury or post-stroke rehabilitation.

\item Elderly Care Facilities \& Nursing Homes: Optimizes institutional care by allowing a small staff footprint to respond instantly to high-priority fall alerts.

\item Smart Home \& IoT Integration: Functions as a localized, network-independent safety module that interfaces seamlessly with residential intranet hubs.

\item Clinical Documentation \& Compliance Monitoring: Automatically logs and formats objective, auditable incident data required for formal medical reports and facility compliance audits.

\item Research \& Development: Serves as an open-hardware reference design and provides a localized, South Asian regional IMU motion dataset for public academic research.
\end{enumerate}